It is believed that the accretion disk around supermassive black holes (BH) at
galactic centres gives rise to UV thermal emission. This emission is
associated with Active Galactic Nuclei (AGNs).

The optical spectra of bright AGNs show additional bright broad emission
lines. Those emission lines arise from the dense gas in the Broad Line Region
(BLR), which is ionized by the UV photons from the accretion disk. See the
sketch to visualise this model.

\ioaafig{2020_DA_1-f1.pdf}

We can assume that the flux of broad emission lines varies in response to the
variation of the UV continuum with a time delay. This time delay should be
proportional to the separation $R_{\mathrm{BLR}}$ between the BH and the BLR.

Assume that the size of the accretion disk is negligible as compared to
$R_{\mathrm{BLR}}$.

\begin{description}
\item[(a)] (1 point) Estimate the time lag (days) between the B-band continuum
  and broad emission line (H-$\beta$) using the light curves shown below. The
  x-axis is in reduced Julian Dates (JD).

  \ioaafig{2020_DA_1-f2.pdf}
\item[(b)] (3 points) Estimate $R_{\mathrm{BLR}}$ in parsecs (pc).
\item[(c)] (2 points) Estimate the angular separation of this region
  $\theta_{\mathrm{BLR}}$ (in arcsec) from the blackhole, if this AGN is
  100\,Mpc away from us.

  It is possible to estimate the mass of the system using the Virial theorem,
  if the velocity dispersion of the gasses in the BLR and the size of the
  system are known. Assume that the masses of the accretion disk and broad
  line region are negligible, as compared to the black hole. The velocity
  dispersion $v_\sigma$ may be estimated from the broadening of the given
  emission line. We will take the corresponding wavelength dispersion to be
  \[ \sigma = \frac{\mathrm{FWHM}}{2.35} \]
  where FWHM is the full width at half maximum of the broad emission line.
\item[(d)] (5 points) Calculate the velocity dispersion $v_\sigma$ in units of
  km\,s$^{-1}$, from the spectral line shown below.

  \ioaafig{2020_DA_1-f3.pdf}
\item[(e)] (4 points) Calculate the mass of the central BH
  ($M_{\mathrm{vir,BH}}$) in a unit of $M_\odot$.
\end{description}
